% SPDX-License-Identifier: Apache-2.0
% covers: Dmem
\datasheet{Dmem}{Vreteno's data memory: a block RAM of words behind an AXI peripheral end.}

\dsfacts{\code{cpu/vreteno/src/dmem.rs}}{\code{vreteno32::dmem::Dmem}}%
{\code{//docs:vreteno}}

\subsection*{Function}

\code{Dmem} is the data memory of the Vreteno machine, a device on its
bus. It holds \code{W} words, \code{DMEM\_WORDS}, 1024, by default,
in four memories of a byte, \code{lane0} to \code{lane3}, one per
lane of the word. A store of a byte or a half writes only its lanes and
reads nothing. The board had a second, \code{Dmem<5, 14, 16384>}, the
stack memory at \code{0x1\_0000} (issue 1278), until the stack window
became a RAM in the core, on its own port (issue 1275).

It is the peripheral client of an \code{AxiPer}, and its link is one
port, \code{bus}, a \code{PerPort} of the four channels: it takes
requests on \code{bus\_req} and write beats on \code{bus\_w}, and
answers on \code{bus\_ans} for a write and \code{bus\_r} for a read.
It answers a read burst of incrementing words a beat a cycle, which
is what the core's instruction cache fills its lines with (issue 1021),
and write bursts, whose beats go to consecutive words and which it
answers once, after the last beat (issue 1208).

\subsection*{Parameters}

\begin{center}\footnotesize
\begin{tabular}{@{}lp{0.7\textwidth}@{}}
\toprule
Parameter & Meaning \\
\midrule
\code{I} & The width of the AXI identifier on \code{bus\_req},
\code{bus\_ans} and \code{bus\_r}, in bits. The tables use 4, as the board
does since its bus has two hosts behind an arbiter (issue 241). \\
\code{AW} & The bits that number the words, 10 by default. \\
\code{W} & The words, \code{1 << AW}, 1024 by default. \\
\bottomrule
\end{tabular}
\end{center}

\subsection*{Address map}

The word a burst names is bits \code{AW + 1} to 2 of its address,
11 to 2 by default. Lane $k$ is
byte $k$ of the word: \code{lane0} is bits 7 to 0, \code{lane3} bits 31
to 24. The memory looks at no other address bits.

On the board, \code{BoardRouter} sends it the page at \code{0x1000}
with mask \code{0xffff\_f000}, so addresses \code{0x1000} to
\code{0x1ffc} are words 0 to 1023. The model's \code{DATA\_BASE} is
the same \code{0x1000}.

\dstables{Dmem}

\subsection*{Behaviour}

\begin{itemize}
\item A write is taken when no other write is held. Its first word
and identifier go into \code{paddr} and \code{pid}, and \code{pend}
is set.
\item The held write's beats are taken as they arrive, the last only
when \code{bus\_ans} has room. Each lane whose strobe bit is set takes
its byte of the beat, and \code{paddr} moves to the next word. The
last beat clears \code{pend}, and \code{bus\_ans} sends \code{OKAY}
with \code{pid}.
\item A read is taken when no write is waiting and the answer register
is free or is being emptied in this cycle. The four lanes are read
into \code{word} at the edge, the synchronous read a block RAM has.
\item The read is answered on \code{bus\_r} the cycle after, with
\code{rid} and \code{OKAY}. A burst's \code{len} goes into
\code{rleft} and the next word's index into \code{raddr}; while
\code{rleft} is not zero, each beat sent reads the next word into
\code{word} and counts \code{rleft} down, and the beat with
\code{rleft} at zero has \code{last} set. A new read is taken only
when the answer register is free or is sending a burst's last beat.
\item While a write is held, no request of either kind is taken.
\item \code{Dmem::with} fills the lanes from bytes before the first
cycle. \code{data\_word} reads a word back for a run or a test.
\end{itemize}

\subsection*{Verification}

\begin{itemize}
\item \code{//cpu/vreteno:lockstep\_test} compares every word of the model's
data memory with the four lanes at the halt, in \code{demo\_program},
\code{a\_multiply\_right\_after\_a\_load} and \code{random\_programs}.
\item \code{//cpu/vreteno:hello\_test}, \code{//cpu/vreteno:hello\_cc\_test}
and \code{//cpu/vreteno:board\_test} put the compiled programs'
constants in it with \code{Dmem::with}, and check what the programs
print.
\item The build simulates \code{Dmem<2>}'s netlist against the trace
of \code{//cpu/vreteno:vreteno} under nvc and Verilator:
\code{//docs:vreteno\_sim\_dmem}\allowbreak\code{\_dmem\_tb\_test} and
\code{//docs:vreteno\_vsim\_dmem\_test}.
\end{itemize}

\subsection*{Used by}

\code{Board}, as \code{dmem} behind the tracker \code{pdmem}. The
demonstration \code{//cpu/vreteno:vreteno}, the lockstep test and
\code{run.rs}. \code{board\_netlist} writes a program's data into
\code{lane0} to \code{lane3} of the board's netlist.

\subsection*{Limits}

It does not read a write burst's length: the last beat ends it. It
checks no address, since the router sends it only its own range, so
an address beyond word \code{W} $-$ 1 wraps. One write is held at a time. Nothing
on the machine fills it at run time except stores, so a program's
constants have to be in it before the first cycle.
\code{cpu/vreteno/board/ax7a200.xdc} constrains the four lanes to block
RAM.
